\begin{table}
    \begin{center}
    % \vspace{-30pt}
        \scalebox{0.85}{
        \begin{tabular}{c | c c c c c c c c c  }
            \Xhline{3\arrayrulewidth} 

            Method & $V^*$ & $V^*_{D.A.}$ & $V^*_{R.P.}$ & MMVP  & IQ-Test & JigSaw    \rule[-1.2ex]{0pt}{4ex}\\
            \Xhline{3\arrayrulewidth} 
             % Qwen2.5-VL & 74.87 & 82.61 & 63.16 & 59 & 69.17 & 28.0 & 50.0 & 41.79	 & 81.82 \\
             % SFT & 75.92	& 85.22	 &61.84	& 62.0 & 63.33	&22.67	&50.67	&41.79	&86.01\\
             \modelname          & 81.7   & 84.4	& 77.6	& 71.7  & 29.3 & 52.0  \\
          \modelname$_{LatentEnd}$ & 39.8   & 32.2 & 51.3  & 19.0  & 6.7  & 13.3  \\
          \modelname$_{MLP Head}$  & 74.4   & 76.5 & 71.1  & 69.7  & 23.3 & 50.0  \\
           \modelname$_{GLU Head}$ & 79.6   & 82.6 & 75.0  & 69.0  & 25.3 & 44.0  \\

         \Xhline{3\arrayrulewidth} 
        \end{tabular}
        } 
        \vspace{-5pt}
        \caption{\small Ablation studies on the 7B model show the standard approach performs best, indicating the LLM natively aligns visual and textual semantics without an extra head. However, the unstable latent end token suggests a need for future work on variable-length reasoning.~}\label{table:ablation}
    \end{center}
\vspace{-15pt}
\end{table}

% We conduct ablation studies on the 7B model by exploring alternative architectural designs and decoding strategies. The standard \modelname configuration achieves the strongest performance, indicating that the LLM can directly accommodate visual semantics for \modelname and textual semantics for next-token prediction without requiring an additional connecting head. We further examine the latent end token as the stopping criterion for the \modelname process. However, its instability during evaluation highlights a promising direction for future research on variable-length latent reasoning. 